% =====================================================================
%  Приложение B. Образцы и их геометрия
%
%  ЖАНР (STYLE.md): учебная литература. Образцы называются физически, по
%  D/P, без внутренних кодовых имён наборов данных. Прослеживаемость —
%  ключи \srckey{B.x} и таблица "образец -> набор" в начале Z_provenance.
% =====================================================================
\section{Образцы и их геометрия}
\label{sec:samples}
\markboth{Приложение. Образцы и их геометрия}{}

Курс опирается на измерения пяти образцов проволочных поляризаторов, различающихся в первую
очередь параметром плотности~$D/P$ (§\ref{sec:electrodynamics}). Это приложение собирает их
геометрию и схему измерения в одном месте~--- разрозненные упоминания по тексту курса (например,
пример~\ref{subsec:no-measurement} §3 или таблица §4) ссылаются на один и тот же набор образцов,
описанных здесь.

% ---------------------------------------------------------------------
\subsection*{Общая схема измерения}

Все измерения выполнены на терагерцовом спектрометре временного разрешения (§0): импульс
известной поляризации проходит через оптическую схему, содержащую исследуемый образец на
поворотном столике, и регистрируется когерентно~--- прибор измеряет само поле, а не его квадрат
(§\ref{subsec:field-vs-intensity}). Схема встречается в проекте в двух конфигурациях:

\begin{itemize}[leftmargin=1.6em,itemsep=0.3em]
  \item \term{Одиночный образец между двумя плёночными поляризаторами.} Исследуемый образец
        стоит между входным и выходным плёночными поляризаторами с высокой заявленной степенью
        экстинкции (в терагерцовом диапазоне~--- практически идеальные проекторы,
        §\ref{subsec:leakage}); поворачивается сам исследуемый образец. Для этой схемы
        одноэлементная модель (§\ref{sec:blanco}--§\ref{sec:drude}) применяется напрямую.
  \item \term{Два реальных элемента тракта.} Оба поляризующих элемента~--- настоящие проволочные
        решётки (не идеальные плёнки); один неподвижен, другой поворачивается. Для этой схемы в
        §1 доказана теорема о выровненном анализаторе, сводящая её (при аккуратной юстировке) к
        той же одноэлементной модели, что и первая схема.
\end{itemize}

Геометрия каждого образца ($D$, $P$) определялась микроскопией; для одного из образцов результат
независимо перепроверен по дифракции видимого света (пример~\ref{subsec:field-vs-intensity} §0,
ключ~0.3)~--- согласие в пределах экспериментальной погрешности. Абсолютная точность определения
диаметра методом микроскопии ограничена оптическим разрешением и способом измерения ширины
отражающей полосы провода; в тексте курса используются значения $D/P$, уже согласованные между
собой и примененные во всех числовых примерах.

% ---------------------------------------------------------------------
\subsection*{Сводная таблица}

\begin{center}
\footnotesize
\setlength{\tabcolsep}{5pt}
\begin{tabular}{@{}p{0.24\linewidth}p{0.09\linewidth}p{0.24\linewidth}p{0.12\linewidth}p{0.24\linewidth}@{}}
\toprule
\textbf{Как назван в курсе} & \textbf{$D/P$} & \textbf{Схема измерения} & \textbf{Число углов} &
\textbf{Где встречается} \\
\midrule
плотная решётка & $0{,}71$ & два реальных элемента тракта (вращается второй; плёнка на входе) &
$19$ & §0-§9, основной пример лестницы моделей (§6) и вырождения (§9) \\
\addlinespace[3pt]
повторная установка (а) & $0{,}69$ & два реальных элемента тракта (вращается первый; без плёнок) &
$19$ & §4 (пример отрицательной энергии тени при равномерных весах) \\
\addlinespace[3pt]
повторная установка (б) & $0{,}69$ & два реальных элемента тракта (вращается первый; обе плёнки) &
$19$ & §4-§5 (таблица офсетов и корреляции) \\
\addlinespace[3pt]
образец средней плотности & $0{,}56$ & одиночный образец между двумя плёнками & $13$ & §2, §3
(модельные расчёты по геометрии) \\
\addlinespace[3pt]
разреженный образец & $0{,}46$ & одиночный образец между двумя плёнками & $25$ & §4-§8 (образец с
почти исчерпанным остатком после опорной модели, ключ~8.1) \\
\addlinespace[3pt]
наиболее разреженный & $0{,}42$ & одиночный образец между двумя плёнками, измерения повторены в
двух независимых сеансах & (не используется напрямую в числовых примерах курса) & §0 (диаметр
пучка, оговорка о конечной ширине) \\
\bottomrule
\end{tabular}
\end{center}
\srckey{B.1}

\begin{warning}[: повторные установки — не независимые образцы]
Три образца с $D/P\approx0{,}69$--$0{,}71$ (плотная решётка и обе повторные установки) физически
представляют собой измерения \textbf{одного и того же} прибора в разных условиях съёмки~---
разные плёнки на входе/выходе, разный из двух элементов вращается,~--- а не три разных образца.
Это уже отмечалось во вводной таблице служебного приложения «Происхождение числовых примеров» и
повторяется здесь: считать их независимыми источниками геометрической информации~--- значит
завысить объём выборки.
\end{warning}

% ---------------------------------------------------------------------
\subsection*{Угловая развёртка: типичная схема}

\begin{intuition}
Угловая развёртка~--- набор измерений при разных положениях поворотного столика, обычно с шагом
$10^\circ$, от отрицательных углов до положительных или до $90^\circ$--$100^\circ$ в одну сторону.
Плотность точек вблизи $90^\circ$ (глубокая тень) для части образцов увеличена~--- именно там,
как показал §1, сосредоточена вся информация о просачивании, и именно там измерение труднее всего
из-за низкого отношения сигнал/шум (§\ref{subsec:mask}).
\end{intuition}

\noindent
Выбор диапазона развёртки~--- не техническая деталь, а прямое проявление содержания §5:
симметричный диапазон даёт хорошо обусловленную матрицу плана, диапазон, зафиксированный только с
одной стороны,~--- обусловленность на два порядка хуже при том же шаге и сопоставимом числе точек.
Образцы курса представляют обе ситуации, и это намеренно: студент, сравнивающий §4-§5 разных
образцов, увидит на настоящих числах, к чему на практике приводит выбор диапазона.

\clearpage
